% ECONOMICS_PROBLEM: {"id": "C72-199", "title": "Improving search-and-mix approximation guarantees", "jel_code": "C72", "source_id": "OP-0132"}
% !TeX program = xelatex
\documentclass[11pt,a4paper]{article}
\usepackage[margin=23mm]{geometry}
\usepackage{fontspec}
\setmainfont{Latin Modern Roman}
\usepackage{amsmath,amssymb,mathtools}
\usepackage{xcolor,enumitem,booktabs,longtable,array,xurl,hyperref}
\definecolor{muted}{HTML}{53636C}
\hypersetup{colorlinks=true,urlcolor=blue,linkcolor=blue}
\setlength{\parindent}{0pt}
\setlength{\parskip}{5pt}
\newcommand{\R}{\mathbb R}
\newcommand{\E}{\mathbb E}
\newcommand{\N}{\mathbb N}
\newcommand{\ind}{\mathbf 1}
\newcommand{\Var}{\operatorname{Var}}
\newcommand{\Cov}{\operatorname{Cov}}
\newcommand{\supp}{\operatorname{supp}}
\DeclareMathOperator*{\argmin}{arg\,min}
\DeclareMathOperator*{\argmax}{arg\,max}
\newcommand{\entrymeta}[3]{{\sffamily\small\color{muted}\textbf{\texttt{#1}}\quad|\quad #2\par #3\par}}
\newcommand{\Problemref}[1]{\texttt{#1}}
\newcommand{\JELref}[2]{\texttt{#2} (JEL \texttt{#1})}
\begin{document}
\section*{Improving search-and-mix approximation guarantees}
\textbf{Problem ID:} \texttt{C72-199}\quad \textbf{JEL:} \texttt{C72}\par
% BEGIN ECONOMICS STATEMENT
\label{problem:OP-0132}
% GAME_THEORY_PROBLEM: {"local_id": "GT-002", "original_number": 2, "kind": "R", "entry_kind": "statement_only", "published": false, "review_required": true, "category": "game_theory"}
\label{game:GT-002}
{\sffamily\small\color{muted}
\textbf{GT-002} | Algorithmic equilibrium and complexity\\
\textbf{R} | Precise algorithm-design target; the source's framework is not an axiomatized complexity class.
\par}
\subsection*{Definitions and mathematical statement}
Let $A,B\in\mathbb Q^{m\times n}\cap[0,1]^{m\times n}$ be explicitly encoded, and let $r_{A,B}$ denote the maximum unilateral payoff gain, as displayed in GT-001. A search-and-mix algorithm first computes finite lists $X=(x^1,\ldots,x^s)\subset\Delta_m$ and $Y=(y^1,\ldots,y^t)\subset\Delta_n$, then outputs
\[
 x=\sum_{a=1}^{s}\lambda_ax^a,\qquad
 y=\sum_{b=1}^{t}\mu_by^b,\qquad
 \lambda\in\Delta_s,\quad\mu\in\Delta_t.
\]
An admissible template specifies the search subroutines, candidate lists, accuracy tolerances, and mixing rule; all must have polynomial bit complexity. The quantitative design target is to exhibit such a fully specified template, an algorithm $\mathcal A$, and a constant $\eta\in(0,1/3)$ satisfying
\[
 \forall A,B,\qquad
 r_{A,B}\bigl(\mathcal A(A,B)\bigr)\leq \frac13-\eta.
\]
For a specified family $\mathfrak S$ of templates, its optimal guarantee can instead be defined by $\alpha(\mathfrak S)=\inf_{\mathcal A\in\mathfrak S}\sup_{A,B}r_{A,B}(\mathcal A(A,B))$.
\subsection*{Short English statement}
Improve the worst-case additive error of a concrete search-and-mix algorithm below one third.
\subsection*{Variants and scope boundaries}
The decomposition alone imposes no meaningful restriction: any algorithm fits it by returning its final answer as a singleton candidate list. Consequently an impossibility theorem for ``the framework'' needs additional restrictions on $\mathfrak S$. The historical benchmark is $1/3+\delta$ for arbitrary $\delta>0$, with runtime polynomial in $1/\delta$, rather than an automatically attained exact $1/3$ guarantee.
\subsection*{Source relationship and status}
[S1] explicitly requests an improvement using search-and-mix methods. The displayed constant-gap target is an editorial mathematical specification of that request. No unverified claim that $1/3$ remains the best bound on the document date is made.
\subsection*{References}
\begingroup\footnotesize\raggedright
{\footnotesize\raggedright
\textbf{[S1]} Hanyu Li, Wenhan Huang, Zhijian Duan, David Henry Mguni, Kun Shao, Jun Wang, and Xiaotie Deng (2023). \emph{A survey on algorithms for Nash equilibria in finite normal-form games}. Locator: Section 3.4, including the DFM22 algorithm; concluding ``Open problems and future directions,'' theoretical item 1. \url{https://arxiv.org/pdf/2312.11063}\par}
\par\endgroup

\subsection*{Common conventions: Source mathematical conventions}

\textbf{Finite games.} For a finite set $A$, $\Delta(A)$ is its probability
simplex. Players choose independent mixed strategies in Nash equilibrium;
correlation is allowed only when explicitly part of the solution concept.
In a game with payoff functions $u_i$ and strategy profile $x$, an additive
$\varepsilon$ Nash equilibrium satisfies
\[
 u_i(x)\geq u_i(x_i',x_{-i})-\varepsilon
 \quad\text{for every player $i$ and every admissible deviation $x_i'$}.
\]
For numerical approximation, payoffs are normalized as locally specified.
An $\varepsilon$ payoff residual is distinct from distance at most
$\varepsilon$ to the set of exact equilibria. Point convergence, convergence
of empirical play, and convergence in distance to an equilibrium set are
also distinct.

\textbf{Computational representations.} Unless stated otherwise, a complexity
question takes rational data with binary encoding; polynomial time is measured
in total bit length. A fixed-accuracy polynomial algorithm is not necessarily
polynomial in $1/\varepsilon$, and neither is necessarily polynomial in
$\log(1/\varepsilon)$. Succinct games and value, demand, or sampling oracles
must declare their interfaces. Exact real solutions require an explicit
representation; an unspecified list of arbitrary real numbers is not a
finite computational output. A claimed algorithmic barrier is meaningful
only for its specified model and reductions.

\textbf{Dynamic games.} Histories, information, observation, and allowable
randomization are part of the model. A stationary strategy depends only on
the current state; a positional strategy in a deterministic graph game
chooses one action at each controlled vertex. Nash equilibrium at an initial
state and equilibrium after every history are different requirements.
An expectation of a pathwise $\liminf$ payoff, a $\liminf$ of expected averages,
a limiting expected average, and a uniform finite-horizon equilibrium are
different criteria. None is silently substituted for another.

\textbf{Allocation and mechanisms.} A complete allocation assigns every item
exactly once unless otherwise stated. Zero-valued goods, negative values,
capacity constraints, feasibility, lotteries, and copies of goods affect
fairness definitions and are specified locally. Dominant-strategy incentive
compatibility, Bayesian incentive compatibility, universal truthfulness,
and truthfulness in expectation impose different inequalities. Whenever
an objective ratio has zero denominator, its convention is stated or those
instances are excluded.

\textbf{Infinite-dimensional models.} Expectations, integrals, stochastic
equations, and optimization objectives require their local regularity,
integrability, filtrations, and admissible domains. A backward--forward
mean-field system is a boundary-value problem; it is not an initial-value
semiflow without an additional construction. Long-time, large-population,
and vanishing-noise limits require an order of limits and a convergence
topology. If a minimum need not be attained, the objective is its infimum
or supremum and the approximation requirement is stated separately.
% END ECONOMICS STATEMENT
\end{document}
